\section{The finite infrastructure inside \(\ipra\)}\label{sec:framework}

\subsection{The object theory actually used}

The ProveIt presentation of \ipra{} is first-order intuitionistic arithmetic
with equality, a symbol and defining equations for every fixed primitive
recursive function, the successor axioms, and induction for quantifier-free
formulas.  The crucial word is \emph{fixed}: we may add the coefficient
function of one series, the verifier of one certificate format, or the step
function of one iteration.  We do not postulate a primitive-recursive universal
evaluator for all primitive-recursive programs.

Every invariant below is represented by an equality of natural-number terms.
For example, a bounded assertion
\[
  \forall i<N\; P(i)
\]
with decidable \(P\) is encoded by the Boolean fold
\[
  \operatorname{All}_P(0)=1,
  \qquad
  \operatorname{All}_P(i+1)=
  \operatorname{All}_P(i)\cdot\chi_P(i).
\]
The assertion is the quantifier-free equation
\(\operatorname{All}_P(N)=1\).  Consequently the ordinary proof ``maintain
this finite invariant for every stage'' is literally quantifier-free
induction.

\subsection{Integers, rationals, and finite algebra}

An integer is coded by a pair \((u,v)\) representing \(u-v\).  Equality is
cross-addition,
\[
 (u,v)=(u',v')\quad\Longleftrightarrow\quad u+v'=u'+v.
\]
A rational is represented by a signed numerator and a positive denominator;
all rational identities are converted to integer identities by multiplying by
the product of the guarded nonzero denominators.  The ProveIt
\code{IPRADerivations.Core.Integers} and \code{Core.Tactic} layers already
implement this difference-ring method and a ring normalizer for object-theory
polynomial identities.

A polynomial of degree at most \(d\) is a finite coefficient vector.  Addition,
multiplication, formal differentiation, evaluation, Euclidean division over a
field, Sylvester determinants and resultants are bounded loops.  A sparse
representation is convenient for the very large universal Tschirnhaus
expressions, but no theorem depends on the choice of representation.

\begin{lemma}[Finite polynomial normalization]\label{lem:poly-normalize}
For each fixed number of variables and each coded degree bound there is a
primitive-recursive normalizer \(\operatorname{nf}\) for integer polynomials
such that \ipra{} proves
\[
 \operatorname{nf}(P)=\operatorname{nf}(Q)
 \quad\Longrightarrow\quad
 P(\mathbf x)=Q(\mathbf x).
\]
The same is true for rational polynomials after denominator clearing.
\end{lemma}

\begin{proof}
Encode a monomial by its exponent vector and a polynomial by a finite list of
coefficient--monomial pairs.  Insertion sort by the lexicographic monomial code,
followed by a bounded fold combining equal monomials, is primitive recursive.
Correctness is proved by induction on the input length.  The invariant is an
equality between the evaluation of the processed prefix and the evaluation of
the normalized accumulator, hence quantifier-free.  Rational coefficients are
handled by a common positive denominator.
\end{proof}

This elementary lemma is the object-theory meaning of every use of
\code{ring}, \code{field\_simp}, or a symbolic resultant in the metatheoretic
Lean and Python artifacts.

\subsection{Regular names for real and complex numbers}

A \emph{regular dyadic name} for a real number is a primitive-recursive
function \(q:\N\to\Q\) satisfying
\begin{equation}\label{eq:regular-name}
  \abs{q(n)-q(m)}\le 2^{-n}+2^{-m}.
\end{equation}
Two names are equal if
\(\abs{q(n)-r(n)}\le 2^{1-n}\) for all \(n\).  A complex name is a pair of real
names.  Arithmetic on names uses explicit error propagation.  On a supplied
rational bound \(\abs{x}\le M\), the polynomial map \(x\mapsto x^5+x\) has an
explicit rational modulus because
\[
 \abs{x^5-y^5}
 \le 5(M+1)^4\abs{x-y}.
\]
after increasing the bound by one to absorb approximation error.

Statements about order use located intervals rather than an illicit global law
of trichotomy for arbitrary reals.  In the Bring construction all branching is
performed on exact rational midpoint evaluations.  The extension from rational
parameters to a named real parameter is then made by the Lipschitz estimate of
Theorem~\ref{thm:bring-lipschitz}; no undecidable sign test on a real number is
required.

\begin{definition}[Effective convergence]
A sequence \(z_n\) converges to \(z\) with modulus \(M\) when
\[
 m,n\ge M(k)\quad\Longrightarrow\quad
 \abs{z_m-z_n}\le 2^{-k}.
\]
every term and \(M\) being primitive recursive.  A power series is asserted to
converge only together with a primitive-recursive tail bound on each named
closed subdisk of its open disk of convergence.
\end{definition}

This is the reading used for the Bring series, hypergeometric functions and
local analytic branches.

\subsection{Truncated power series}

For a commutative coded ring \(R\), define
\[
  T_N(R)=R[X]/(X^{N+1})
\]
as vectors \((a_0,\ldots,a_N)\).  Cauchy multiplication is
\[
  (ab)_n=\sum_{i=0}^{n}a_i b_{n-i}.
\]
Composition \(A(B)\) is a bounded sum whenever \(B_0=0\).  If \(U_0=1\), the
inverse of \(U\) in \(T_N(R)\) is produced by the finite geometric series in
\(1-U\), or recursively coefficient by coefficient.  All these operations are
primitive recursive in \(N\).

The finite Lagrange theorem used later can be stated without completed series.

\begin{theorem}[Finite Lagrange inversion, reusable form]\label{thm:finite-lagrange}
Let \(\Phi\in T_N(\Q)\).  Starting from \(G_0=0\), define
\[
 G_{r+1}=X\Phi(G_r)\pmod{X^{N+1}}.
\]
Then \(G_{N+1}\) is the unique vector with zero constant term satisfying
\(G=X\Phi(G)\).  For \(1\le n\le N\),
\begin{equation}\label{eq:finite-lagrange}
 n[X^n]G=[W^{n-1}]\Phi(W)^n.
\end{equation}
Every operation and the proof of the theorem are primitive recursive in \(N\),
and the theorem is derivable in \ipra{} by quantifier-free induction.
\end{theorem}

\begin{proof}
If two vectors agree below degree \(m\), their images under
\(A\mapsto X\Phi(A)\) agree below degree \(m+1\).  Thus the successive iterates
stabilize one coefficient per step, which gives existence and uniqueness after
\(N+1\) steps.  Formula~\eqref{eq:finite-lagrange} follows from the finite
Laurent-residue substitution identity
\[
 \Res Q(F(W))F'(W)=\Res Q(W).
\]
after truncating every Laurent vector to the finite exponent interval that can
contribute to the residue.  The proof on Laurent monomials is a telescoping
coefficient identity; linearity gives the general case.  This is exactly the
architecture of the existing ProveIt Lagrange-inversion packet.
\end{proof}

Writing the equation with the factor \(n\) on the left avoids division in the
object theory.  Division by \(n>0\) is used only in the rational metanotation.

\subsection{Finite fields, permutations, and groups}

For the Galois arguments we use two complementary encodings.

\begin{enumerate}
\item For a concrete polynomial modulo a concrete prime, elements of
      \(\F_p[X]/(f)\) are coefficient vectors of length \(\deg f\).  Addition,
      multiplication, gcd, Frobenius powers and Bézout certificates are
      bounded algorithms.
\item An abstract finite group is a natural number \(n\) together with an
      \(n\times n\) multiplication table and tables for identity and inverse.
      The group axioms, subgroup closure, commutators, derived series, orbits
      and permutation actions are bounded Boolean folds.
\end{enumerate}

A permutation of five letters is a five-entry vector.  There are only
\(5!=120\) candidates.  The executable certificate enumerates all 156 actual
subgroups of \(S_5\), identifies the 20 transitive ones, and checks that the 18
solvable transitive subgroups have orders \(5,10\), or \(20\) and are contained
in a conjugate of \(\operatorname{AGL}(1,5)\).  This is a finite computation,
so its trace can be reproduced by a primitive-recursive verifier inside
\ipra.

\begin{lemma}[Finite certificate reflection]\label{lem:certificate-reflection}
Let \(V(c)\in\{0,1\}\) be a fixed primitive-recursive verifier, and suppose
\ipra{} proves the uniform soundness implication
\[
 V(c)=1\longrightarrow R(c).
\]
after all bounded assertions in \(R\) have been Boolean-folded.  Then every
explicit accepted certificate \(\bar c\) yields an \ipra{} proof of
\(R(\bar c)\).
\end{lemma}

\begin{proof}
The closed computation \(V(\bar c)=1\) reduces by the defining equations of the
primitive-recursive symbols.  Modus ponens with the uniform soundness theorem
gives \(R(\bar c)\).  No reflection principle for arbitrary proofs is being
assumed; this is only evaluation of one fixed decidable verifier.
\end{proof}

\subsection{Reuse of the checked ProveIt library}

The following dependencies are already available in the repository.

\begin{itemize}
\item \code{IPRADerivations.Core.Semiring}, \code{Core.Integers},
      \code{Core.Order}, \code{Core.Powers}, \code{Core.Search} and
      \code{Core.Pairing} provide the basic arithmetic, bounded search and
      coding infrastructure.
\item \code{IPRADerivations.Algebra.PolynomialFormulas.quadratic},
      \code{cardano}, and \code{ferrari} are closed derivations of the exact
      factorization identities needed to solve the auxiliary equations of
      degrees two, three and four.
\item The transfer theorem \code{pra\_iff\_native} transports these ordinary
      \ipra{} derivations to the repository's very small native scalar core.
\end{itemize}

The present Lean blueprint is designed to add new object-theory modules above
this checked base rather than to replace it with Mathlib field reasoning.
